\documentclass[11pt,a4paper]{article}
\usepackage[margin=2.5cm]{geometry}
\usepackage[T1]{fontenc}
\usepackage{booktabs,longtable,array}
\usepackage{listings}
\usepackage{xcolor}
\usepackage{hyperref}
\usepackage{enumitem}

\definecolor{codebg}{gray}{0.95}
\lstset{basicstyle=\ttfamily\footnotesize,backgroundcolor=\color{codebg},breaklines=true,
        breakatwhitespace=false,frame=single,columns=fullflexible,keepspaces=true,showstringspaces=false}
% Red placeholders: search for TODO in this file and replace them with YOUR real results.
\newcommand{\todo}[1]{\textcolor{red}{\textbf{[#1]}}}
\newcommand{\repourl}{\url{https://github.com/YOUR-USERNAME/cryptography-network-security-exam}}

\title{Cryptography and Network Security\\ \large Integrated Situation --- Technical Report\\
       \normalsize ETTCS801 \,|\, ULK Polytechnic Institute}
\author{\todo{Your full name} \\ \todo{Registration number}}
\date{\today}

\begin{document}
\maketitle
\tableofcontents
\bigskip

%---------------------------------------------------------------------
\section{Introduction}
A polytechnic institute stores student records on a central server and transfers files between two campuses.
A security review found weak staff passwords, outdated software, unencrypted file transfers and guest-network
access to the records server, and repeated connection attempts from an unfamiliar external address were observed.
This report documents the risk assessment, the Python encryption and integrity toolkit, and the firewall
rules produced to address these findings. All source files are in the GitHub repository (Section~\ref{sec:conclusion}).

%---------------------------------------------------------------------
\section{Risk Assessment}
\subsection{Assets, vulnerabilities and consequences}
\begin{longtable}{@{}p{0.3cm}p{3.3cm}p{4.2cm}p{6.3cm}@{}}
\toprule
\# & Asset & Vulnerability & Possible consequence\\ \midrule\endhead
R1 & Staff accounts and credentials & Weak passwords; repeated external login attempts &
Account takeover; records read, altered or deleted; misuse under a staff identity.\\[4pt]
R2 & Student records server & Guest network can reach the server; outdated software &
Scanning and exploitation of known flaws; data breach, ransomware, outage; reputational and legal harm.\\[4pt]
R3 & Record files in transit between campuses & Unencrypted file transfer &
Eavesdropping or man-in-the-middle; records disclosed or silently modified.\\
\bottomrule
\end{longtable}

\subsection{Risk ranking}
Likelihood and impact are scored from 1 (low) to 5 (high); the score is their product.
\begin{center}
\begin{tabular}{@{}clccc@{}}
\toprule
Rank & Risk & Likelihood & Impact & Score\\ \midrule
1 & R1 -- Weak staff passwords & 5 & 5 & 25\\
2 & R2 -- Guest access / outdated software & 4 & 5 & 20\\
3 & R3 -- Unencrypted transfers & 3 & 4 & 12\\ \bottomrule
\end{tabular}
\end{center}
\textbf{Reasoning.} R1 ranks first because attacks are already occurring and weak passwords need little skill to
defeat; one compromised account exposes all data the user can reach. R2 is second: guests are many and untrusted, and
known exploits exist for outdated software, with very high impact because all records reside on one server. R3 ranks
third because interception requires a position on the inter-campus path, although success exposes and allows tampering with complete files.

\subsection{Recommended controls}
\begin{center}
\begin{tabular}{@{}lp{11cm}@{}}
\toprule
Risk & Control\\ \midrule
R1 & Strong password policy (minimum 12 characters, no reuse), multi-factor authentication, account lockout / rate limiting.\\
R2 & Network segmentation with firewall rules denying guest access to the server, plus a regular patching schedule.\\
R3 & Encrypt files before transfer (AES-256-GCM) and use secure channels (SFTP/TLS/VPN), with SHA-256 integrity checks.\\
\bottomrule
\end{tabular}
\end{center}

%---------------------------------------------------------------------
\section{Encryption and Integrity Application}
\subsection{Design}
The programme \texttt{src/secure\_records.py} is written in Python and uses the \texttt{cryptography} library.
\begin{itemize}[nosep]
  \item \textbf{Encryption} -- AES-256-GCM authenticated encryption. A random 256-bit key is generated once and stored
        \emph{outside} the repository (default \texttt{\textasciitilde/.ulk\_keys/records.key}); the tool refuses a key path inside a git
        repository. Each encryption uses a fresh random 96-bit nonce. The output file is
        \texttt{header (4 bytes) $\|$ nonce (12 bytes) $\|$ ciphertext $\|$ GCM tag (16 bytes)}.
  \item \textbf{Decryption} -- reads the nonce, decrypts and authenticates. A wrong key or any modified byte raises an
        authentication failure, reported as a clean error message. With \texttt{--original} the decrypted file's SHA-256
        is compared with the original's to prove the contents match.
  \item \textbf{Integrity} -- \texttt{hash} stores the SHA-256 digest of a file in a baseline file (\texttt{hashes.json});
        \texttt{check} recomputes the digest and reports \emph{UNCHANGED} or \emph{MODIFIED}.
  \item \textbf{Error handling} -- missing files, directories, empty files, invalid or missing keys, corrupt or non-encrypted
        input and corrupt baselines are caught and reported with a message and exit code, never a traceback.
\end{itemize}

\subsection{Test evidence}
Ten automated unit tests (round trip, ciphertext differs from plaintext, tamper detection, wrong key, missing key,
missing/empty/invalid input, hash change detection, CLI error code) all pass:
\begin{lstlisting}
$ python3 -m unittest discover -s tests -v
...
Ran 10 tests in 0.008s
OK
\end{lstlisting}
Full demonstration output (\texttt{tests/run\_demo.sh}); the two SHA-256 values match after decryption, and the
modified file is flagged:
{\footnotesize
\lstinputlisting[basicstyle=\ttfamily\tiny]{../tests/crypto_demo_output.txt}}

%---------------------------------------------------------------------
\section{Network Traffic Filtering}
\subsection{Laboratory setup}
\todo{Describe your lab: server IP, guest network, staff network, other/external client, firewall tool (iptables), and the service the assessor specified.}

\subsection{Firewall rules}
Rules are in \texttt{firewall/records\_server\_firewall.sh} and use a dedicated iptables chain so they can be removed cleanly.
In order of evaluation:
\begin{enumerate}[nosep]
  \item Accept established/related traffic (and loopback) so existing sessions keep working.
  \item \textbf{Block guest $\rightarrow$ server (all ports):} \texttt{-s GUEST\_NET -j DROP} (logged, rate limited).
  \item \textbf{Permit staff $\rightarrow$ named service:} \texttt{-s STAFF\_NET -p tcp --dport PORT -j ACCEPT}.
  \item \textbf{Block all other inbound access to the service:} \texttt{-p tcp --dport PORT -j DROP}.
\end{enumerate}
\begin{lstlisting}
iptables -A RECORDS_FILTER -m conntrack --ctstate ESTABLISHED,RELATED -j ACCEPT
iptables -A RECORDS_FILTER -s $GUEST_NET -j DROP
iptables -A RECORDS_FILTER -s $STAFF_NET -p tcp --dport $SERVICE_PORT \
         -m conntrack --ctstate NEW -j ACCEPT
iptables -A RECORDS_FILTER -p tcp --dport $SERVICE_PORT -j DROP
iptables -A RECORDS_FILTER -j RETURN
iptables -I INPUT 1 -j RECORDS_FILTER
\end{lstlisting}
\todo{Replace variables with your real values if you want the report to show them.}

\subsection{Test procedure and results}
One permitted and two blocked connections were tested with \texttt{nc -zv -w 5 SERVER\_IP PORT}.
\begin{center}
\begin{tabular}{@{}llp{3.3cm}p{3.6cm}@{}}
\toprule
Test & Source & Expected & Actual\\ \midrule
T1 permitted & Staff client & Connection succeeds & \todo{paste result}\\
T2 blocked & Guest client & Timeout / dropped & \todo{paste result}\\
T3 blocked & Other client & Timeout / dropped & \todo{paste result}\\ \bottomrule
\end{tabular}
\end{center}
\todo{Paste the terminal output or a screenshot from filter\_tests.md, and the iptables packet counters.}

%---------------------------------------------------------------------
\section{Conclusion}\label{sec:conclusion}
The assessment ranked weak passwords, guest access to the records server, and unencrypted transfers as the main risks.
The toolkit addresses the last two directly: authenticated AES-256-GCM encryption protects files in transit and at rest, SHA-256
baselines detect tampering, and the firewall removes the guest path to the server while limiting the service to the staff
network. Weak passwords are addressed by policy and MFA, which are organisational controls outside the code. Limitations: the key must
be distributed securely between campuses (a key-management system would be needed at scale), and firewall rules cover only
the tested service. Future work: automated patching, centralised logging and intrusion detection.

\medskip
\noindent\textbf{GitHub repository:} \repourl

\end{document}
